\section{线性结构的教材级展开}

线性数据结构看起来简单，实际最容易被学薄。数组、字符串、链表、栈、队列和双端队列不是几组 API 名称，而是不同访问模式的承诺：数组承诺按下标常数访问，链表承诺已知节点附近的局部重连，栈承诺只处理最近未解决对象，队列承诺按进入顺序扩展，双端队列承诺两端候选都能快速进出。算法题里的“优化”常常就是把原本昂贵的访问模式换成某个结构天然支持的访问模式。经典教材会反复强调这一点，因为如果只会背容器名字，遇到变体题时就无法判断模板为什么失效。

学习线性结构时可以用三条线索串起来。第一条是物理存储：元素是否连续，是否会整体搬迁，节点地址是否稳定，是否能被 CPU 预取。第二条是抽象接口：能不能随机访问，能不能从中间删除，能不能只看一端，能不能同时维护两端。第三条是算法不变量：慢指针左侧是什么，栈内保存什么未解决对象，队列当前层表示什么距离，窗口左端为什么可以丢弃。真正掌握数据结构，就是能把这三条线索同时讲清楚。

\begin{principlebox}{访问模式先于容器名称}
选择容器前先问访问模式：按下标频繁访问，优先连续数组；只在尾部追加且要排序或 DP，优先 \code{vector}；需要头尾都弹入弹出，考虑 \code{deque}；需要已知节点处常数移动，考虑链表；需要最近未匹配结构，用栈；需要按层扩展，用队列。容器不是越复杂越好，而是越贴合访问模式越好。
\end{principlebox}

\topic{数组和动态数组：连续内存、下标语义和原地维护}

数组的核心能力是“位置稳定且可计算”。第 i 个元素的位置可以由起始地址加上 i 乘以元素大小得到，所以按下标访问是常数时间。这一条性质让数组天然适合二分、前缀和、差分数组、排序、双指针、动态规划表和邻接表。数组的弱点也来自连续性：在中间插入或删除一个元素，后面的元素必须整体搬移；容量不足时，动态数组要申请新内存并移动已有元素。

\code{std::vector} 的概念模型可以理解成三个指针：起始位置、当前结束位置、容量结束位置。\code{size()} 来自当前结束位置减起始位置，\code{capacity()} 来自容量结束位置减起始位置。尾部追加时如果容量够，只在尾部构造一个元素；容量不够时，申请更大内存、移动旧元素、释放旧内存。这个过程解释了两个常见问题：为什么 \code{push_back} 的均摊复杂度是常数，为什么扩容后旧指针、引用、迭代器会失效。

数组算法要特别重视下标区间语义。左闭右闭区间 \code{[left, right]} 适合描述题目原始边界，左闭右开区间 \code{[left, right)} 更适合循环和切片，因为长度正好是 \code{right - left}，空区间表示为 \code{left == right}。前缀和通常使用左闭右开语义：\code{prefix[i]} 表示前 i 个元素之和，区间 \code{[left, right)} 的和就是 \code{prefix[right] - prefix[left]}。这种设计把下标零的特殊情况变成统一公式，是教材里非常重要的哨兵思想。

原地维护类题目本质是在数组上划分语义区间。删除有序数组重复项中，慢指针左侧是已经写好的唯一值前缀；颜色分类中，数组被划成零区、一号区、未知区和二区；移动零中，慢指针左侧是已经保留相对顺序的非零元素；下一个排列中，后缀是最长非递增区，枢轴和后缀重排决定字典序的下一步。题目不同，但证明方式相同：每一步循环只扩大已经正确的区域，未知区域严格缩小，被排除或交换出去的元素不会再破坏不变量。

先看颜色分类。题目只有三种颜色，最直接的暴力办法是排序，或者先扫描三遍分别统计 0、1、2 的数量，再按数量覆盖回数组。排序当然能得到结果，但它没有利用值域只有三类这个条件；计数覆盖虽然是线性的，却把题目变成“先统计再重写”，没有训练原地分区。进一步观察，扫描过程中真正需要维护的不是完整频次，而是四段区间：左边已经确定为 0，中间已经确定为 1，右边已经确定为 2，剩下的是未知区。于是优化自然变成荷兰国旗式的三指针维护，而不是直接背一段交换代码。

\begin{principlebox}{图示：颜色分类的四段区间}
\begin{tabular}{@{}p{0.31\linewidth}p{0.31\linewidth}p{0.26\linewidth}@{}}
区间 & 语义 & 变化方式\\
\code{[0, next_red)} & 已经确定为 0 & 只向右扩大\\
\code{[next_red, current)} & 已经确定为 1 & 只向右扩大\\
\code{[current, next_blue]} & 尚未检查 & 每轮缩小\\
\code{(next_blue, n)} & 已经确定为 2 & 只向左扩大\\
\end{tabular}
\end{principlebox}

\begin{lstlisting}[language=C++,caption={颜色分类：四段区间不变量}]
void sort_colors(std::vector<int>& nums)
{
    constexpr int RED = 0;
    constexpr int WHITE = 1;
    constexpr int BLUE = 2;

    int next_red = 0;
    int current = 0;
    int next_blue = static_cast<int>(nums.size()) - 1;

    while (current <= next_blue) {
        if (nums[current] == RED) {
            std::swap(nums[next_red], nums[current]);
            ++next_red;
            ++current;
        } else if (nums[current] == BLUE) {
            std::swap(nums[current], nums[next_blue]);
            --next_blue;
        } else {
            ++current;
        }
    }
}
\end{lstlisting}

这段代码的关键不是三指针本身，而是区间解释：\code{[0, next_red)} 全是 0，\code{[next_red, current)} 全是 1，\code{[current, next_blue]} 未知，\code{(next_blue, n)} 全是 2。遇到 2 时不能推进 \code{current}，因为从右侧换回来的元素还没有检查；遇到 0 时可以推进，因为换回来的位置属于已经扫描过的一号区或就是当前。这个细节正是很多错误实现的来源。

数组变种还包括稀疏数组、位图、环形缓冲区和压缩坐标数组。位图把布尔集合压到二进制位，适合值域有限且只需要存在性的问题；环形缓冲区用固定数组和头尾下标模拟队列，适合实时系统和滑动窗口；离散化把大值域映射到连续小编号，适合树状数组、线段树和频次统计。刷题时常见的“坐标压缩”不是数学技巧，而是把无法直接开数组的值域转成数组能支持的下标访问。

\topic{字符串：字符序列、编码假设和子串成本}

字符串是数组的特例，但题目里的字符串又比数组多几个坑。第一，字符集决定计数结构：小写英文可以用 26 长度数组，ASCII 可以用 128 或 256 长度数组，Unicode 则不能简单按字节计数。第二，\code{substr} 会构造新字符串，热循环中频繁使用会引入复制、分配和哈希成本。第三，字符串算法常常需要区分子串和子序列：子串必须连续，适合窗口、哈希、KMP 和中心扩展；子序列不要求连续，常用动态规划或贪心扫描。

固定长度窗口题的本质是相邻窗口只差一个进入字符和一个离开字符。找到所有字母异位词、排列包含、固定长度最大元音数都属于这一类。变长窗口题则依赖合法性可以通过移动左端恢复，例如最小覆盖子串中右端扩展增加覆盖能力，满足后移动左端尝试变短。若窗口状态不会随左右端单调变化，例如数组含负数且目标是区间和，就不能硬套滑动窗口。

字符串匹配类题目有多条路线。暴力匹配每个起点再逐字符比较，最坏可能 \complexity{O(nm)}；KMP 通过前缀函数避免主串指针回退；滚动哈希把子串比较转成哈希值比较，但要处理冲突；Trie 把多个模式串的公共前缀合并，适合单词搜索 II、前缀统计和自动补全。面试时不一定要写出所有高级算法，但必须能说明它们解决的是哪种重复工作。

排列包含是固定长度窗口的典型例子。题目问 \code{text} 中是否存在某个长度等于 \code{pattern.size()} 的子串，字符频次和 \code{pattern} 完全相同。暴力做法是枚举每个起点，截出长度为 m 的子串，再排序比较或重新统计频次。相邻两个候选子串只差左边移出的一个字符和右边进入的一个字符，如果每次从零统计，就把 m 个字符的绝大部分重复计算了。优化点因此很明确：窗口长度固定，右端每前进一步，只增量加入一个字符，并在窗口超过长度时移出一个字符；窗口计数等于需求计数时返回真。

\begin{principlebox}{图示：固定长度窗口的增量更新}
\begin{tabular}{@{}p{0.2\linewidth}p{0.72\linewidth}@{}}
旧窗口 & \code{text[left ... right]} 中已经维护好 26 个字符计数。\\
右端进入 & 新字符 \code{text[right + 1]} 的计数加一。\\
左端离开 & 若窗口长度超过 \code{pattern.size()}，旧字符 \code{text[left]} 的计数减一。\\
判断答案 & 窗口长度固定后，比较 \code{window == need}，不用重新统计整段。\\
\end{tabular}
\end{principlebox}

\begin{lstlisting}[language=C++,caption={字符串窗口：用数组计数表达字符集假设}]
bool contains_permutation(const std::string& text, const std::string& pattern)
{
    constexpr int ALPHABET_SIZE = 26;
    if (text.size() < pattern.size()) {
        return false;
    }

    std::array<int, ALPHABET_SIZE> need{};
    std::array<int, ALPHABET_SIZE> window{};
    for (const char ch : pattern) {
        ++need[ch - 'a'];
    }

    for (std::size_t right = 0; right < text.size(); ++right) {
        ++window[text[right] - 'a'];
        if (right >= pattern.size()) {
            --window[text[right - pattern.size()] - 'a'];
        }
        if (window == need) {
            return true;
        }
    }
    return false;
}
\end{lstlisting}

这段代码必须附带条件：输入只含小写英文字母。如果题目没有这个条件，数组下标 \code{ch - 'a'} 就没有意义。经典教材会反复强调“前提条件”，因为算法复杂度和正确性都依赖前提。把字符集换成 Unicode 后，代码结构仍然可以是窗口，但计数结构需要换成哈希表或更明确的编码处理。

\topic{链表：节点身份、局部重连和哨兵节点}

链表的价值不在于“删除一定比数组快”，而在于节点身份稳定和局部重连便宜。已知一个节点或它的前驱时，插入删除只需要改少数指针；但如果要寻找第 k 个节点，仍然必须从头走 k 步。链表题考的是结构不变量：哪些节点已经接好，哪些节点还没有处理，当前节点的后继是否已经保存，操作后是否仍然能到达未处理部分。

单链表、双链表、循环链表和侵入式链表各有不同用途。单链表空间少，适合反转、合并、快慢指针；双链表能从节点向前向后移动，适合 LRU、删除已知节点和双端队列；循环链表适合轮转和约瑟夫类问题；侵入式链表把链接字段放进业务对象本身，Linux 内核大量使用这种方式来减少额外分配并精确控制对象生命周期。刷题通常使用平台给出的 \code{ListNode}，但理解这些变体能帮助你解释为什么 LRU 需要哈希表加双向链表。

哨兵节点是链表题的统一边界工具。删除头节点、合并空链表、分隔链表、反转前 k 个节点时，如果直接操作真实头节点，经常要写一堆特殊分支；引入虚拟头节点后，所有操作都变成“修改某个前驱的 \code{next}”。这和前缀和多开一个 \code{prefix[0]} 是同一种思想：用一个额外状态把边界并入普通逻辑。

两两交换节点如果从数组视角出发，暴力办法是把节点值复制到数组里，交换相邻值后再写回链表。这样能得到相同值序列，却没有真正交换节点，也无法迁移到要求保持节点身份的题。若坚持在链表上做，问题可以缩小到一个局部：前驱节点 \code{previous} 后面有两个待交换节点 \code{first} 和 \code{second}，再后面是下一段入口 \code{after}。优化不是减少时间复杂度，因为暴力重建也是线性；优化的是结构质量和可迁移性：只改三条边，保持已经处理好的前缀不动，未处理后缀仍然可达。

\begin{principlebox}{图示：两两交换的局部重连}
\begin{tabular}{@{}p{0.18\linewidth}p{0.74\linewidth}@{}}
交换前 & \code{previous -> first -> second -> after}\\
第一条边 & \code{previous -> second}，让前缀接到交换后的头。\\
第二条边 & \code{second -> first}，完成这一对内部交换。\\
第三条边 & \code{first -> after}，把未处理后缀接回链表。\\
推进位置 & \code{previous = first}，因为 \code{first} 已成为已处理前缀的尾部。\\
\end{tabular}
\end{principlebox}

\begin{lstlisting}[language=C++,caption={两两交换链表节点：哨兵和局部三指针}]
ListNode* swap_pairs(ListNode* head)
{
    ListNode dummy;
    dummy.next = head;

    ListNode* previous = &dummy;
    while (previous->next != nullptr && previous->next->next != nullptr) {
        ListNode* first = previous->next;
        ListNode* second = first->next;
        ListNode* after = second->next;

        previous->next = second;
        second->next = first;
        first->next = after;

        previous = first;
    }
    return dummy.next;
}
\end{lstlisting}

这段代码的局部不变量是：\code{previous} 之前的链表已经交换完成，\code{previous->next} 是当前待交换对的第一个节点，\code{after} 保存下一段入口。先保存 \code{after}，再改三条边，最后把 \code{previous} 移到这一对交换后的尾部。K 个一组反转也是同一模式，只是局部段更长，需要先检查是否有 k 个节点，再在段内反转并接回。

链表常见变种包括环检测、环入口、相交链表、回文链表、链表排序、复制带随机指针链表和 LRU 缓存。环检测用快慢指针，核心是速度差；相交链表比较的是节点地址，不是节点值；回文链表需要找到中点、反转后半段、比较并最好恢复结构；链表排序适合自底向上归并，避免递归栈；随机指针复制需要原节点到新节点的映射，或者把复制节点插入原链表再拆分。每个变种都不是新模板，而是围绕“节点身份和链接关系”展开。

\topic{栈和单调栈：最近关系、表达式和候选淘汰}

栈保存的是最近尚未解决的结构。括号匹配中，最近的左括号必须先匹配；路径简化中，最近进入的目录可能被 \code{..} 取消；表达式求值中，最近的运算符和操作数形成局部归约；DFS 中，栈保存待展开路径。普通栈的正确性来自后进先出，单调栈的正确性来自候选淘汰。

单调栈不是为了“保持栈有序”而有序，而是为了让新元素到来时能结算一批旧元素。每日温度、下一个更大元素、柱状图最大矩形、接雨水都可以用这个模型解释。栈里保存尚未找到右边界的下标；新元素若破坏单调性，就说明某些栈顶元素的右边界已经确定。被弹出的元素不会再回来，因此总复杂度是摊还线性。

表达式类题目则体现栈的另一面：保存上下文。逆波兰表达式已经把优先级编码到顺序里，只需要数字栈；基本计算器 II 需要遇到乘除立即结合、加减延迟求和；带括号的计算器需要在进入括号时保存外层结果和符号。写表达式题时，先定义 token 类型、当前数字、当前运算符、栈中保存内容，再动手编码，错误会少很多。

逆波兰表达式可以先按题意暴力模拟：从左到右读 token，遇到数字暂存，遇到运算符就回头找离它最近的两个尚未参与归约的数字，计算后再把结果放回序列。这个做法的瓶颈是“最近尚未归约对象”的查找和删除很别扭，因为序列中间会不断被替换。观察运算符的定义后会发现，真正需要的永远是最近的两个操作数，且先出现的操作数在左，后出现的操作数在右。这正是栈的职责：数字入栈，运算符弹出右操作数和左操作数，归约结果再入栈。

\begin{principlebox}{图示：逆波兰表达式的栈归约}
\begin{tabular}{@{}p{0.2\linewidth}p{0.72\linewidth}@{}}
读到数字 & 压入栈，表示一个尚未被运算符消费的操作数。\\
读到运算符 & 弹出 \code{rhs}，再弹出 \code{lhs}，顺序不能反。\\
完成归约 & 计算 \code{lhs op rhs}，把结果重新压栈。\\
最终状态 & 所有 token 扫完后，栈顶就是整个表达式的值。\\
\end{tabular}
\end{principlebox}

\begin{lstlisting}[language=C++,caption={逆波兰表达式：栈中保存待归约操作数}]
int eval_rpn(const std::vector<std::string>& tokens)
{
    std::vector<int> stack;
    stack.reserve(tokens.size());

    for (const std::string& token : tokens) {
        if (token == "+" || token == "-" || token == "*" || token == "/") {
            const int rhs = stack.back();
            stack.pop_back();
            const int lhs = stack.back();
            stack.pop_back();

            if (token == "+") {
                stack.push_back(lhs + rhs);
            } else if (token == "-") {
                stack.push_back(lhs - rhs);
            } else if (token == "*") {
                stack.push_back(lhs * rhs);
            } else {
                stack.push_back(lhs / rhs);
            }
        } else {
            stack.push_back(std::stoi(token));
        }
    }
    return stack.back();
}
\end{lstlisting}

这里弹出的第一个数是右操作数，第二个数是左操作数。减法和除法不能交换顺序，这是表达式栈题最常见的细节。工程实现中，如果 token 很多且转换成本重要，可以先写一个更明确的解析函数；但刷题训练重点是栈状态语义。

\topic{队列、双端队列和环形缓冲区}

队列保存的是时间顺序或层次顺序。BFS 第一次到达某个状态即为最短距离，前提是所有边权相等且访问标记在入队时设置。多源 BFS 只是把多个起点同时作为第零层，腐烂橘子、地图中离最近零的距离、从边界反向标记安全区域都可以这样处理。队列不负责“选最优候选”，它只负责“按层推进”；如果边权不同，就需要堆或其他最短路结构。

双端队列比队列多了从两端进出的能力，因此能维护窗口候选。单调队列保存的不是窗口内所有元素，而是仍可能成为答案的候选下标。队头是当前答案，队尾负责删除被新元素支配的旧候选。滑动窗口最大值、和至少为 K 的最短子数组、受限子序列和都依赖类似的支配关系。

环形缓冲区是队列的底层变体。固定数组配合头尾下标，可以避免频繁分配；当尾部走到数组末尾时回到开头。它适合实现固定容量队列、日志缓冲、生产者消费者队列和某些实时系统。刷题中不常要求手写环形队列，但理解它能帮助你明白为什么 \code{deque} 和 \code{queue} 的接口强调两端操作，而不是随机插入。

\begin{lstlisting}[language=C++,caption={固定容量环形队列：头尾下标表达状态}]
class CircularQueue {
public:
    explicit CircularQueue(int capacity)
        : values_(static_cast<std::size_t>(capacity)), capacity_(capacity)
    {
    }

    bool push(int value)
    {
        if (this->size_ == this->capacity_) {
            return false;
        }
        this->values_[static_cast<std::size_t>(this->tail_)] = value;
        this->tail_ = (this->tail_ + 1) % this->capacity_;
        ++this->size_;
        return true;
    }

    std::optional<int> pop()
    {
        if (this->size_ == 0) {
            return std::nullopt;
        }
        const int value = this->values_[static_cast<std::size_t>(this->head_)];
        this->head_ = (this->head_ + 1) % this->capacity_;
        --this->size_;
        return value;
    }

private:
    std::vector<int> values_;
    int capacity_ = 0;
    int head_ = 0;
    int tail_ = 0;
    int size_ = 0;
};
\end{lstlisting}

这段代码保留 \code{size_}，是为了区分空队列和满队列。如果只用 \code{head_ == tail_}，空和满都会满足这个条件，需要浪费一个槽位或额外标记。真实工程中还要考虑并发、异常安全和容量为零的输入约束；刷题版本则要把状态不变量讲清楚：\code{head_} 指向下一个弹出位置，\code{tail_} 指向下一个写入位置，\code{size_} 表示当前元素数。

\topic{线性结构的实际应用和变种题谱}

数组的实际应用包括排序缓冲、DP 表、图邻接表、前缀统计、差分更新、坐标压缩后的频次数组。字符串应用包括日志解析、词典匹配、搜索建议、DNA 序列重复检测、表达式解析。链表应用包括 LRU、内核对象链、空闲块管理、任务队列和需要稳定节点地址的缓存。栈应用包括函数调用、解析器、撤销操作、表达式求值、单调候选。队列应用包括 BFS、任务调度、消息缓冲、多源扩散。双端队列应用包括滑动窗口极值、0-1 BFS、双端搜索和单调 DP 优化。

变种题复盘可以按结构归类。数组原地维护：26、27、75、88、189、283；前缀与差分：303、304、560、1109、1094；字符串窗口：3、76、438、567；链表重连：19、21、24、25、92、206；链表身份：141、142、160、138；普通栈：20、71、150、224、227、394；单调栈：42、84、496、503、739；队列 BFS：102、199、200、994、1091；单调队列：239、862、1438、1425。每组题都要问同一个问题：结构保存什么，什么时候加入，什么时候移除，移除是否会丢答案。

\begin{labbox}
线性结构实验：第一，用 \code{vector}、\code{deque}、\code{list} 分别顺序扫描一百万个整数，比较实际时间并解释缓存局部性。第二，把字符串窗口题写成 \code{substr} 加排序和计数窗口两版，记录输入长度增长时的差异。第三，把链表反转的每一步打印为节点地址，确认比较的是节点身份。第四，把单调栈中的每个下标入栈出栈次数记录下来，验证摊还复杂度。第五，手写环形队列并测试空、满、绕回和连续弹入弹出。
\end{labbox}

\topic{线性结构变种题谱}

数组变种可以分成七类。第一类是原地过滤，例如删除重复项、移除元素、移动零、稳定划分。这类题的共同点是输出仍放在原数组前缀，慢指针表示下一个写入位置，快指针负责扫描未知区。第二类是双向收缩，例如盛水容器、两数之和 II、回文判断。这类题依赖左右边界比较能排除一侧。第三类是前后缀贡献，例如除自身以外数组乘积、接雨水的前后缀解法、最短无序连续子数组。第四类是矩阵边界模拟，例如螺旋矩阵、旋转图像、矩阵置零。第五类是差分和扫描线，例如航班预订、拼车、区间加法。第六类是坐标压缩后用数组维护频次或排名。第七类是环形数组，例如轮转数组、循环队列、下一个更大元素 II。

这些变种应该统一回到区间不变量。原地过滤要说明写入区不会被未来读操作破坏；双向收缩要说明被移走的一侧不可能产生更优答案；前后缀贡献要说明当前位置答案由左右两边已经缓存的信息合成；矩阵模拟要说明四条边界何时收缩，剩一行或一列时不会重复访问；差分数组要说明每个区间更新如何转化为两个端点事件。题目名称不同，但证明方式不是散的。

字符串变种可以按“连续、非连续、多模式”分。连续子串题包括无重复最长子串、最小覆盖子串、异位词窗口、回文子串、重复 DNA 序列。非连续子序列题包括判断子序列、最长公共子序列、不同的子序列。多模式题包括单词搜索 II、替换单词、搜索建议系统。连续子串常用窗口、哈希、中心扩展和 KMP；子序列常用双指针或 DP；多模式常用 Trie。复盘时先判断目标是否要求连续，能直接避免很多错误方向。

链表变种可以按“位置、结构、身份、排序、缓存”分。位置类题包括删除倒数第 N 个、旋转链表、分隔链表，核心是快慢指针和前驱节点。结构类题包括反转链表、两两交换、K 个一组反转、重排链表，核心是保存后继和局部接回。身份类题包括环形链表、环入口、相交链表，核心是比较节点地址。排序类题包括合并两个有序链表、合并 K 个链表、排序链表，核心是节点重连和归并。缓存类题包括 LRU 和 LFU，核心是哈希定位和链表维护顺序。

栈变种可以按“匹配、归约、单调、上下文”分。匹配类题包括有效括号、删除无效括号、最长有效括号；归约类题包括逆波兰表达式、基本计算器、字符串解码；单调类题包括每日温度、下一个更大元素、柱状图最大矩形、接雨水；上下文类题包括路径简化、文件系统解析、嵌套列表解析。判断一个栈题属于哪一类，关键看栈里保存的是符号、操作数、下标、局部字符串，还是上一层上下文。

队列变种可以按“层次、双端、单调、优先级”分。普通队列服务 BFS 层次，双端队列服务两端扩展和 0-1 BFS，单调队列服务窗口极值和 DP 转移最大值，优先级队列服务动态极值。不要把这些队列混成一个概念。普通 BFS 队列不按权重选点，所以不能处理不同边权；单调队列不是为了先进先出，而是为了删除被支配候选；优先级队列不是按时间顺序，而是按优先级顺序。

\begin{principlebox}{线性结构复盘表}
每做一道线性结构题，都在题解末尾补一张小表：输入是否连续，是否允许修改，是否需要保持相对顺序，是否要比较节点身份，窗口是否合法依赖什么状态，候选什么时候过期，元素最多进入和离开结构几次。能填完这张表，说明你掌握的是结构，不是题号。
\end{principlebox}

\topic{线性结构的底层成本模型}

经典教材不会只给出一张复杂度表，因为复杂度表把太多真实代价藏起来了。线性结构的底层成本至少有五个维度：随机访问、顺序扫描、局部插入删除、内存分配、缓存局部性。数组和 \code{vector} 在随机访问和顺序扫描上很强，因为第 i 个元素的位置可以直接计算；链表在已知节点附近重连很强，但每走一步都要先加载下一个节点地址；\code{deque} 介于两者之间，分段连续让两端扩张便宜，却让纯顺序扫描的局部性通常弱于 \code{vector}。

这五个维度会直接改变刷题方案。最长递增子序列的 \complexity{O(n\log n)} 解法需要频繁在一个有序尾数组上二分，尾数组必须支持下标访问，所以 \code{vector} 合适；LRU 缓存需要把任意已知节点移到表头并删除表尾，\code{list} 合适；滑动窗口最大值需要从两端删除候选，\code{deque} 合适；环形日志缓冲不需要无限扩容，只需要固定容量覆盖旧值，数组加头尾下标更合适。容器不是按照“高级程度”选择，而是按照操作组合选择。

再看插入删除的说法。说链表删除是 \complexity{O(1)}，必须补上前提：已经有待删节点的前驱或迭代器。如果你只有一个值，还要从头找这个值，查找仍然是 \complexity{O(n)}。说 \code{vector} 尾部插入是均摊 \complexity{O(1)}，也必须补上前提：它可能偶尔扩容，扩容时旧地址全部失效。说 \code{deque} 两端插入是 \complexity{O(1)}，也不等于它适合所有随机访问密集任务。高质量答案要把这些前提说出来。

\begin{deepdive}
CPU 缓存解释了为什么很多教材和工程代码偏爱连续数组。处理器从内存读取一个整数时，通常会把相邻的一小段字节一起载入缓存；如果下一次访问正好落在这段缓存行里，就不必再等主内存。数组顺序扫描天然吃到这个收益。链表节点可能散落在堆上，访问下一个节点前必须先读出当前节点里的指针，硬件很难提前知道下一步在哪里。两段代码同为 \complexity{O(n)}，实际速度可能差出很多，这不是大 O 错了，而是大 O 的抽象层次不同。
\end{deepdive}

\topic{数组模型：从区间到事件}

数组题要先判断自己维护的是“区间状态”还是“位置事件”。区间状态题关心某个连续范围内的和、最大值、字符频次、不同元素数；位置事件题关心某个下标开始或结束了一段影响。前缀和把区间查询转成两个端点的差，差分数组把区间更新转成两个端点事件。二者是同一思想的两个方向：前缀和适合多次查询，差分适合多次更新后统一还原。

差分数组特别适合讲清“事件化”的思维。航班预订、拼车、区间加法这类题，如果对每个区间逐个位置累加，复杂度会被区间长度拖垮。差分数组只在左端点加一个事件，在右端点后一位减一个事件；最后扫描一次，把事件累积成每个位置的真实值。它把“影响一整段”改写成“从这里开始多一个影响，从那里之后少一个影响”。

\begin{lstlisting}[language=C++,caption={差分数组：区间更新变成端点事件}]
std::vector<int> apply_range_additions(
    int length,
    const std::vector<std::array<int, 3>>& operations)
{
    std::vector<int> difference(static_cast<std::size_t>(length + 1), 0);

    for (const auto& operation : operations) {
        const int left = operation[0];
        const int right = operation[1];
        const int delta = operation[2];
        difference[static_cast<std::size_t>(left)] += delta;
        if (right + 1 < length) {
            difference[static_cast<std::size_t>(right + 1)] -= delta;
        }
    }

    std::vector<int> values(static_cast<std::size_t>(length), 0);
    int current = 0;
    for (int i = 0; i < length; ++i) {
        current += difference[static_cast<std::size_t>(i)];
        values[static_cast<std::size_t>(i)] = current;
    }
    return values;
}
\end{lstlisting}

这段代码的关键是端点语义。若操作作用在闭区间 \code{[left, right]}，就在 \code{left} 处增加影响，在 \code{right + 1} 处撤销影响。若题目使用半开区间 \code{[left, right)}，撤销位置就是 \code{right}。很多差分题错在闭区间和半开区间混用。写题解时先定义区间语义，再写公式，代码才不会靠样例猜。

数组还能表达“排名”和“值域”。计数排序、桶排序、位图、坐标压缩都利用这一点。若值域很小，直接用值作下标，查询和统计都可以很快；若值域很大但实际出现的值不多，就把出现过的值排序后映射到 \code{0..m-1}，这就是离散化。树状数组统计逆序对、区间和个数、动态排名时，经常先离散化，因为原始数值可能达到十亿，不能直接开数组。

\topic{字符串模型：窗口、自动机和字典}

字符串题的第一分叉是连续还是不连续。连续子串题可以用窗口、哈希、KMP、Manacher、后缀结构；不连续子序列题更多走双指针和动态规划；多个模式串同时匹配时，要考虑 Trie 或 AC 自动机。只要这个分叉判断错，后面的模板就容易全错。比如最小覆盖子串必须连续，适合滑动窗口；最长公共子序列不要求连续，窗口没有意义；单词搜索 II 同时查很多单词，逐个单词 DFS 会重复大量前缀，Trie 才能合并重复工作。

滑动窗口的本质是“窗口合法性可以增量维护”。无重复最长子串中，合法性由每个字符出现次数是否超过一决定；最小覆盖子串中，合法性由所有必需字符是否都达到需求决定；替换后的最长重复字符中，合法性由窗口长度减最高频字符数是否不超过替换次数决定。这些题的状态不同，但都能随着左右端移动做常数或近似常数更新。

\begin{principlebox}{窗口题的三句证明}
第一，右端扩张时，窗口状态如何更新；第二，窗口不合法时，左端如何移动并恢复合法；第三，左端移动不会错过答案的原因是什么。正数数组求最短和至少为 K 的子数组能用窗口，是因为和随右端扩张不减；数组允许负数时，这条单调性消失，就要换成前缀和加单调队列。
\end{principlebox}

字符串工程应用还要重视编码和生命周期。刷题题面常说只含小写字母，此时 \code{array<int, 26>} 既快又清晰；日志、路径、URL、自然语言文本则不能这样写。若函数只临时观察一段字符串，\code{string_view} 能减少复制；若要把子串作为长期 key 存进哈希表，就必须拥有一份真正的 \code{string}。算法题为了清晰常省略这些工程前提，但教材应该把它们说清楚。

\topic{链表模型：缓存、空闲链和节点生命周期}

链表题表面是指针题，底层是节点生命周期题。数组的元素位置由下标决定，链表的元素位置由指针关系决定；数组移动元素会改变位置，链表重连节点不会改变节点身份。这个差异让链表适合缓存、内存分配器、内核对象队列和需要稳定节点地址的系统结构。Linux 内核常用侵入式链表：链表节点字段嵌入业务对象中，通过节点地址反推出外层对象。这样做能减少额外分配，也让对象生命周期由系统自己精确控制。

刷题里最能体现工程味的是 LRU。它需要两个操作同时快：按 key 找到缓存项，按新旧顺序删除最久未使用项。哈希表能快速定位 key，却不知道谁最旧；链表能维护顺序，却不能按 key 快速定位。把二者组合起来，哈希表保存 key 到链表迭代器，链表保存从新到旧的缓存项，访问时用 \code{splice} 把节点移到表头。

如果直接写暴力 LRU，可以用一个数组或链表保存缓存项。每次 \code{get} 时线性查找 key，找到后把它移到最前；每次 \code{put} 时也先线性查找，容量满了删除最后一个。这个版本容易理解，但每个操作最坏都是 \complexity{O(n)}，慢点集中在“按 key 定位”这一步。换成单独哈希表可以把定位降到常数，却失去新旧顺序；换成单独链表能维护顺序，却仍然无法快速定位 key。于是组合结构是从两个瓶颈推出来的：哈希表负责定位，双向链表负责顺序，二者通过迭代器连接。

\begin{principlebox}{图示：LRU 的两个索引}
\begin{tabular}{@{}p{0.2\linewidth}p{0.72\linewidth}@{}}
哈希表 & \code{key -> list iterator}，负责常数时间定位缓存项。\\
双向链表 & \code{newest <-> ... <-> oldest}，负责维护访问新旧顺序。\\
\code{get} & 哈希定位节点，再把节点移动到链表头部。\\
\code{put} & 新 key 插入头部；容量满时删除链表尾部，并同步删除哈希表 key。\\
\end{tabular}
\end{principlebox}

\begin{lstlisting}[language=C++,caption={LRU 缓存：哈希定位和链表维护顺序}]
class LruCache {
public:
    explicit LruCache(int capacity) : capacity_(capacity) {}

    int get(int key)
    {
        const auto found = this->position_by_key_.find(key);
        if (found == this->position_by_key_.end()) {
            return -1;
        }
        this->items_.splice(this->items_.begin(), this->items_, found->second);
        return found->second->value;
    }

    void put(int key, int value)
    {
        if (this->capacity_ == 0) {
            return;
        }

        const auto found = this->position_by_key_.find(key);
        if (found != this->position_by_key_.end()) {
            found->second->value = value;
            this->items_.splice(this->items_.begin(), this->items_, found->second);
            return;
        }

        if (static_cast<int>(this->items_.size()) == this->capacity_) {
            const int expired_key = this->items_.back().key;
            this->position_by_key_.erase(expired_key);
            this->items_.pop_back();
        }

        this->items_.push_front(Entry{key, value});
        this->position_by_key_[key] = this->items_.begin();
    }

private:
    struct Entry {
        int key = 0;
        int value = 0;
    };

    int capacity_ = 0;
    std::list<Entry> items_;
    std::unordered_map<int, std::list<Entry>::iterator> position_by_key_;
};
\end{lstlisting}

这个实现的核心不在代码长度，而在职责分离。\code{position_by_key_} 回答“这个 key 在哪里”，\code{items_} 回答“谁最新，谁最旧”。\code{splice} 只改链表链接，不复制节点，所以哈希表里保存的迭代器仍然指向同一个节点。如果用 \code{vector} 存缓存顺序，把元素移到开头会导致大量搬移，也会破坏保存的位置。这个题是数据结构组合的典型教材案例。

链表变体题也要从节点生命周期解释。复制带随机指针链表时，原节点和新节点之间必须建立映射，否则随机边无法接回；回文链表如果反转后半段，比较后最好恢复原结构，避免修改输入的副作用；排序链表适合自底向上归并，因为链表不支持随机访问，快速排序的分区优势不明显；K 个一组反转需要先确认剩余节点够 K 个，再局部反转，否则尾部不足一组的节点不能被破坏。这些细节都来自“节点关系是结构本身”。

\topic{栈、队列和调度模型}

栈和队列不只是简单容器，它们分别表达两类调度策略。栈表达最近上下文优先处理，适合解析嵌套结构、撤销操作、显式 DFS 和单调候选；队列表达先到状态优先处理，适合 BFS、任务缓冲和层次传播；双端队列表达两端都有特殊含义，适合单调窗口、0-1 BFS 和双向搜索。

0-1 BFS 是双端队列的代表。边权只有 0 和 1 时，普通 Dijkstra 用堆当然能做，但双端队列更贴合结构：走 0 权边不增加距离，应该放到队头尽快处理；走 1 权边距离加一，放到队尾。队列里状态的距离仍然近似按非降顺序展开，因此可以得到最短路。

从暴力角度看，带权最短路可以枚举路径，但路径数量会指数级增长；普通 BFS 会把每条边都当成同样代价，遇到 0 权边和 1 权边混在一起时，层数不再等于距离。Dijkstra 的堆版本是通用优化，但在 0/1 边权下，堆顶排序只需要区分“不增加距离”和“距离加一”两类。于是双端队列成为更轻的结构：0 权转移放队头，1 权转移放队尾，保证更小距离的状态先被继续扩展。

\begin{principlebox}{图示：0-1 BFS 的双端队列}
\begin{tabular}{@{}p{0.2\linewidth}p{0.72\linewidth}@{}}
队头方向 & 0 权边不增加距离，新状态放到队头，尽快继续扩展。\\
队尾方向 & 1 权边让距离加一，新状态放到队尾，排在当前距离层之后。\\
核心不变量 & 队列中的候选距离近似非降，先处理的状态不会被更小距离状态越过。\\
失效条件 & 边权不是 0 或 1 时，双端队列不再足够，应改用 Dijkstra。\\
\end{tabular}
\end{principlebox}

\begin{lstlisting}[language=C++,caption={0-1 BFS：双端队列表达两类边权}]
std::vector<int> zero_one_bfs(
    const std::vector<std::vector<std::pair<int, int>>>& graph,
    int source)
{
    constexpr int INF = 1'000'000'000;
    std::vector<int> distance(graph.size(), INF);
    std::deque<int> deque;

    distance[static_cast<std::size_t>(source)] = 0;
    deque.push_front(source);

    while (!deque.empty()) {
        const int node = deque.front();
        deque.pop_front();

        for (const auto& [next, weight] : graph[static_cast<std::size_t>(node)]) {
            const int candidate = distance[static_cast<std::size_t>(node)] + weight;
            if (candidate >= distance[static_cast<std::size_t>(next)]) {
                continue;
            }
            distance[static_cast<std::size_t>(next)] = candidate;
            if (weight == 0) {
                deque.push_front(next);
            } else {
                deque.push_back(next);
            }
        }
    }
    return distance;
}
\end{lstlisting}

这段代码比普通 BFS 多了一个判断：边权为 0 的状态放前面，边权为 1 的状态放后面。它也比堆版 Dijkstra 少了对数因子。成立条件非常窄：边权必须只有 0 和 1；如果出现权重 2、5 或任意非负数，就该回到 Dijkstra。教材式学习要把这种边界讲清楚，否则学生会把一个漂亮技巧误用到不满足条件的图上。

\topic{线性结构题目的讲解标准}

一题讲清楚，至少要包含五层。第一层是暴力方法：枚举什么，复杂度为什么高。第二层是瓶颈：重复扫描、重复计算、候选过多、边界分支太多，还是查询能力不足。第三层是结构选择：为什么数组、链表、栈、队列或双端队列能降低瓶颈。第四层是不变量：哪些区间、节点、候选或层次已经正确。第五层是边界条件：空输入、单元素、重复值、负数、环、容量为零、字符集假设和是否允许修改输入。

以滑动窗口最大值为例，暴力方法是每个窗口扫描 K 个元素；瓶颈是相邻窗口大量重叠，却重新求最大；结构选择是双端队列，因为窗口左端会过期，右端会加入新候选；不变量是队列下标从前到后递增，值从前到后单调不增，队头始终在窗口内；边界条件是 K 非法、数组为空、重复元素是否用小于还是小于等于淘汰。只有这些都说清楚，题解才像教材，而不是模板注释。

以反转链表为例，暴力方法并不是重点，因为它本身已经是线性；重点在结构不变量。循环开始时，\code{previous} 指向已经反转好的前缀头，\code{current} 指向尚未处理部分的第一个节点，\code{current} 之后的链仍保持原方向。每次先保存 \code{next}，再把当前节点接到已反转前缀前，最后推进两个指针。边界条件是空链表和单节点链表。能用这套语言解释，两两交换和 K 个一组反转就不再是新题。

最后要把线性结构放进复盘表里。数组题复盘区间语义和是否原地修改；字符串题复盘字符集、连续性和子串复制成本；链表题复盘节点身份、前驱和后继保存顺序；栈题复盘保存的是符号、下标、操作数还是上下文；队列题复盘层次、距离和访问标记时机；双端队列题复盘候选淘汰和过期规则。这样复盘一段时间后，你看到新题会先想到访问模式，而不是先搜索记忆里的题号。
